\section{Потоки и разрезы}

\subsection{Домашнее задание}
\begin{enumerate}

  \item
    \begin{enumerate}
      \item
        Разбейте вершины ориентированного графа на циклы. Т.е. каждая вершина должна быть покрыта ровно одним циклом.
        Либо скажите, что это невозможно. $\O(E V)$.

        \textbf{Решение:}

        \begin{itemize}
          \item Разбиение вершин на непересекающиеся циклы существует тогда и только тогда,
          когда можно выбрать подмножество рёбер, в котором у каждой вершины ровно одна входящая
          и ровно одна исходящая дуга (граф является объединением вершинно-непересекающихся циклов).

          \item Построим двудольный граф $H = (L \cup R,\; E_H)$:
          \begin{itemize}
            \item $L = \{u^+ \mid u \in V\}$ --- ``исходящая``'' копия каждой вершины,
            \item $R = \{v^- \mid v \in V\}$ --- ``входящая``'' копия каждой вершины,
            \item для каждого ребра $(u \to v)$ в исходном графе добавляем ребро $(u^+,\; v^-)$ в $H$.
          \end{itemize}

          \item Паросочетание в $H$, насыщающее вершину $u^+ \in L$, задаёт выбор
          единственной исходящей дуги из $u$; паросочетание, насыщающее $v^- \in R$, задаёт выбор
          единственной входящей дуги в $v$. Поэтому \emph{совершенное паросочетание} (насыщающее
          все $2V$ вершин $H$) в точности соответствует подграфу с $\deg^+ = \deg^- = 1$ у каждой
          вершины, т.е.\ разбиению на циклы.
          Если совершенного паросочетания нет --- разбиение невозможно.

          \item Сложность: Ищем максимальное паросочетание в $H$ методом
          увеличивающих путей: запускаем DFS из каждой из $V$ вершин левой доли,
          каждый DFS обходит не более $O(E)$ рёбер графа $H$ (рёбра $H$ биективны рёбрам
          исходного графа). Итого $O(VE)$.
          В конце проверяем, что паросочетание совершенно.

          \item Восстановление циклов: По паросочетанию для каждого $u$ известен
          единственный <<следующий>> сосед. Обходя цепочки соседей, выписываем циклы за $O(V)$.
        \end{itemize}

      \item
        Разбейте вершины ориентированного ациклического графа на минимальное количество путей. $\O(E V)$.

        \textbf{Решение:}

        \begin{itemize}
          \item Построим тот же двудольный граф $H = (L \cup R,\; E_H)$, что и в пункте (a):
          $L = \{u^+\}$, $R = \{v^-\}$, рёбра соответствуют рёбрам исходного DAG.

          \item Каждое ребро паросочетания $(u^+, v^-)$
          означает: путь, проходящий через $u$, продолжается в $v$.
          Вершина $u^+$, не насыщенная паросочетанием, является \emph{началом} некоторого пути;
          вершина $v^-$, не насыщенная паросочетанием, является \emph{концом} некоторого пути.
          Путь --- это связная цепочка из ребёр паросочетания.
          Поскольку граф ациклический, цепочки конечны и не зацикливаются.

          \item При максимальном паросочетании размера $M$ число незасыщенных
          вершин в $L$ равно $V - M$ (каждая такая вершина начинает свой путь).
          Значит, минимальное число путей $= V - M$.

          \item Любое разбиение на пути задаёт паросочетание в $H$
          (для каждого внутреннего ребра пути $u \to v$ включаем $(u^+, v^-)$ в паросочетание).
          Число путей при этом равно $V$ минус размер этого паросочетания.
          Максимизируя паросочетание, минимизируем число путей.

          \item Сложность: Поиск максимального паросочетания --- $V$ запусков DFS по
          $O(E)$ каждый, итого $O(VE)$.

          \item Восстановление разбиения: Для каждой незасыщенной вершины $u^+ \in L$
          начинаем новый путь и жадно идём по ребрам паросочетания: если $(u^+, v^-)$ ---
          ребро паросочетания, следующая вершина пути --- $v$. Повторяем, пока цепочка не обрывается.
        \end{itemize}
    \end{enumerate}

  \item
    Дан граф и выделенные вершины $s$, $t$. Нужно проверить, правда ли существует единственный минимальный $s$-$t$ разрез.
    \begin{enumerate}
      \item $\O(poly(V, E))$

        \textbf{Решение:}

        \begin{itemize}
          \item Найдём максимальный поток любым алгоритмом (например, алгоритмом Диница за $O(V^2 E)$). Построим остаточный граф $G_f$.

          \item Введём два специальных множества:
          \begin{itemize}
            \item $A$ --- вершины, \emph{достижимые} из $s$ в $G_f$ (обход в глубину/ширину из $s$).
            \item $B$ --- вершины, из которых $t$ \emph{недостижим} в $G_f$, т.е.\ $B = V \setminus \{v \mid t \text{ достижим из } v \text{ в } G_f\}$ (обход в глубину/ширину по обратным рёбрам из $t$, взять дополнение).
          \end{itemize}

          \item $(A, V \setminus A)$ --- наименьший (по $S$-стороне) минимальный разрез; $(B, V \setminus B)$ --- наибольший. Всегда $A \subseteq B$. Все минимальные разрезы --- это в точности разрезы вида $(C, V \setminus C)$, где $A \subseteq C \subseteq B$.

          \item Следовательно, минимальный разрез единственен тогда и только тогда, когда $A = B$.

          \item Проверка: вычисляем $A$ и $B$ двумя обходами графа $G_f$ за $O(V + E)$ и сравниваем.
          Суммарная сложность определяется поиском максимального потока: $O(V^2 E)$ или $O(poly(V, E))$.
        \end{itemize}

      \item $\O(E)$ при условии, что нам уже известен максимальный поток (сам поток, не только его величина).

        \textbf{Решение:}

        \begin{itemize}
          \item Так как нам дан сам поток, остаточный граф $G_f$ можно построить за $O(E)$:
          для каждого ребра $(u \to v)$ с пропускной способностью $c$ и потоком $f$:
          \begin{itemize}
            \item если $f < c$: добавляем прямое остаточное ребро $(u \to v)$ с остатком $c - f$,
            \item если $f > 0$: добавляем обратное остаточное ребро $(v \to u)$ с остатком $f$.
          \end{itemize}

          \item Вычисляем $A$: обход (DFS/BFS) из $s$ по $G_f$ за $O(V + E)$.

          \item Вычисляем $B$: обход (DFS/BFS) из $t$ по рёбрам $G_f$, взятым в обратном направлении, за $O(V + E)$; берём дополнение.

          \item Отвечаем: единственный разрез $\Leftrightarrow$ $A = B$. Проверка равенства двух множеств за $O(V)$.

          \item Итоговая сложность: $O(E)$ (построение $G_f$ + два обхода + сравнение).
        \end{itemize}
    \end{enumerate}

  \item \exceptgroup{Мишунина} \\
    Есть заказы и инструменты. Для каждого заказа известен список инструментов, который
    нужен, чтобы его выполнить.
    Каждый инструмент сделан умелыми японскими рабочими, поэтому бесконечно прочный,
    его можно один раз купить и много раз использовать.
    У каждого инструмента есть цена $p_i$.
    У каждого заказа есть прибыль, которую можно получить, выполнив заказ.
    Вы -- бедный китайский рабочий. У вас изначально нет инструментов, но зато вы можете под нулевой процент в банке взять
    сколь угодно большой кредит, чтобы купить инструментов.
    \begin{enumerate}
      \item
        Вопрос: какую максимальную прибыль вы можете получить?

      \item
        А теперь тот же вопрос, но ещё есть разные скидочные предложения! \\
        Скидка позволяет два инструмента $i, j$ купить по специальной цене $d \colon \max(p_i, p_j) < d < p_i + p_j$.
        Каждый инструмент присутствует не более чем в одном скидочном предложении.
    \end{enumerate}

  \item
    В неориентированном графе без кратных рёбер необходимо удалить минимальное число рёбер так,
    чтобы увеличилось количество компонент связности.
    \begin{enumerate}
      \item $\O(V \cdot Flow)$, где $Flow$ --- время работы какого-нибудь алгоритма поиска максимального потока.
      \item $\O(E^2)$.
    \end{enumerate}

  \item \onlygroup{Мишунина} \\
    Пусть целевая функция $f : \F_2^n \rightarrow \mathbb{R}$ определена как:
    $$f(x_0, x_1, \dots x_{n-1}) = \sum_{i = 0}^{n - 1} w_i(x_i) + \sum_{i < j} w_{ij}(x_i, x_j)$$
    При этом для каждой $w_{ij}$ выполняется $w_{ij}(0, 0) + w_{ij}(1, 1) \le w_{ij}(0, 1) + w_{ij}(1, 0)$.
    Сведите минимизацию $f$ к минимальному разрезу в графе полиномиального размера с неотрицательными весами. \\
    \textit{Подсказка: попробуйте решить задачу отдельно для двух слагаемых и представить:}
    $$ w_{ij}(x_i, x_j) = A_{ij} + B_{ij} x_i + C_{ij} x_j + D_{ij} \mathbb{I}[ x_i \ne x_j ] $$

\subsection*{Дополнительные задачи}

  \item
    Пусть алгоритм Форда-Фалкерсона ищет дополняющие пути \texttt{DFS}-ом с фиксированным порядком
    просмотра рёбер в каждой вершине, и пропускает максимум по найденному дополняющему пути.
    Постройте вместе с порядком ребер (или докажите, что такого не существует) ориентированный граф:
    \begin{enumerate}
      \item
        с целочисленными пропускными способностями, на котором алгоритм работает за экспоненту от $V$
        (пропускные способности могут быть большими, но должны быть полиномиальной суммарной длины).

      \item
        с вещественными пропускными способностями, на котором алгоритм не завершается за конечное количество шагов.

      \item
        с вещественными пропускными способностями, на котором алгоритм не сходится к корректному значению
        величины максимального потока.
    \end{enumerate}

  \item
    Какое максимальное количество уголков можно разместить на шахматной доске $n \times m$ с дырками?
    Уголком называется фигура, состоящая из трех клеток: центральная клетка черного цвета и две соседних
    с ней белых клетки со смежными сторонам.

\end{enumerate}
